\subsection{The remaining Ihara characters: no new arithmetic
nonvanishing theorem}
\label{sec:kt-ihara-gate}

The carried K17-GREENBERG-KERNEL states that finiteness of
$X^+$ implies $\Omega_p^*=\Omega_p$. Its source is
K17.tex, Theorem K17-GREENBERG-KERNEL, beginning at
line~9678. The supplied proof framework when odd torsion
persists is K17-CUBIC-ODD, beginning at line~10303.
Its conclusion requires, in addition to the structural
freeness and cyclicity assumptions, Pair-NV and Triple-NV
for the \emph{actual} arithmetic companion errors.
In the notation of that theorem, Triple-NV is
\[
 \left.
 \mathcal C_{o;a,b,c}(T_1,T_2,T_3)
 \right|_{(1+T_1)(1+T_2)(1+T_3)=\alpha}\ne0
                                                               \tag{IK.1}
\]
for at least one matching ordered odd triple for each root
$\alpha$ of the relevant odd companion polynomial.
The condition is (CU.1), K17.tex lines~10429--10440.

No arithmetic proof of (IK.1), and no positive-rank
replacement for it, has been obtained in this workstream.
The known finite-level $K$-group vanishings cannot be
silently replaced by all-weight control of the entire
plus Iwasawa module at an arbitrary prime. Brown's rational
mixed-Tate freeness likewise does not prove this integral
arithmetic nonvanishing condition. These are distinct
inputs in the supplied reduction, not alternative wordings
of the same theorem.

There is, however, an exact conditional dependence on the
main question. If all $K_{4k}(\mathbf Z)$ vanish, the carried
Kurihara equivalence gives Vandiver at every prime. The
carried K17-VANDIVER-KERNEL (K17.tex, beginning at line~9305)
then gives $\Omega_p^*=\Omega_p$ for all odd primes; K17-P2-MOMENT
already supplies $p=2$. Thus item~5 would follow from a
completed item~1, without a separate cubic calculation.
No converse and no new unconditional prime are asserted.

In the language of the \emph{Esquisse}, \S2, this remains
a question about the arithmetic action on the pro-$p$
level of the Teichm\"uller tower; the rational unipotent
mixed-Tate quotient does not by itself determine its
pro-$p$ kernel. The supplied LM6 A.44--A.45 and K17
Q3-IRREG rows therefore remain open precisely at the
arithmetic gate (IK.1), unless the main all-weight
vanishing theorem is supplied.

